\documentclass[11pt]{article}

% ---------------------------------------------------------------------------
% What actually helps on-device RAG at 6 GB
% Makata AI Edge Lab — consolidated study of PoCs 001-020
% Build:  latexmk -pdf main.tex   (or: pdflatex; bibtex; pdflatex x2)
% License: CC BY 4.0 (matches the lab's write-ups)
% ---------------------------------------------------------------------------

\usepackage[margin=1in]{geometry}
\usepackage[T1]{fontenc}
\usepackage[utf8]{inputenc}
\usepackage{booktabs}
\usepackage{array}
\usepackage{amsmath}
\usepackage{textcomp}
\usepackage{microtype}
\usepackage[numbers,sort&compress]{natbib}
\usepackage[hidelinks]{hyperref}
\usepackage{xcolor}

\newcommand{\poc}[1]{PoC~#1}
\newcommand{\repo}{\href{https://github.com/ralph-mattew/makata-edge-lab}{makata-edge-lab}}

\title{\textbf{What actually helps on-device RAG at 6\,GB:\\
a pre-registered study of retrieval, query formulation,\\
and the answer step for contract QA with a 1.5B model}}

\author{Ralph Mattew Palomaria\\
  Makata AI Edge Lab\\
  \texttt{ralphmattew@gmail.com} \quad ORCID: \texttt{0000-0000-0000-0000}\\
  \small\href{https://makata.ai}{makata.ai} \,\textbar\, \repo}

\date{Preprint --- \today}

\begin{document}
\maketitle

\begin{abstract}
On-device retrieval-augmented generation (RAG) on a 6\,GB phone is constrained in ways
cloud RAG is not: a sub-2B generator, a few hundred megabytes for the embedder, a short
context window, and thermal limits. We report a pre-registered study of twenty experiments
(\poc{001}--\poc{020}) that test, one mechanism at a time, what raises answer quality for
document question-answering under those limits. The target stack is the one shipped in a
document-chat app: Qwen2.5-1.5B (Q4) as the single on-device model, EmbeddingGemma-300M for
retrieval, and llama.cpp; the benchmark is CUAD contract QA. Each experiment fixes its
question, hypotheses and decision rule before the first run, publishes raw outputs, and is
labelled MEASURED/OBSERVED/INFERRED. We find (i) a strong embedder is the single largest
retrieval lever (recall@3 35\%$\rightarrow$61\%), while BM25-only, model-generated chunk
context, extra retrieval depth, and clause-aware chunking do not help on this benchmark;
(ii) the residual ``hard categories'' are a \emph{query-vocabulary} problem, fixed cheaply by
deterministic keyword expansion (Exclusivity recall 7\%$\rightarrow$32\%) rather than by
scoring, depth, or chunking; and (iii) once the right text is retrieved, a 1.5B generator is
the bottleneck for clause-style answers---so a generator-free extractive answer step
outperforms it (hard-category span-F1 15.7$\rightarrow$20.7) at lower RAM. Several widely-cited
techniques did \emph{not} transfer to this regime, and we report those negative results in
full. All code, raw results and per-experiment commit hashes are public and open to
replication.
\end{abstract}

\section{Introduction}
Retrieval-augmented generation is usually studied with large models and ample compute.
Shipping it \emph{on a phone} changes the constraints: the generator is a quantized sub-2B
model, the embedder must fit in a few hundred megabytes, the context window is short, and
sustained inference throttles under heat. Techniques that help in the cloud may not survive
this regime, and the only way to know is to measure on the target stack.

This paper consolidates twenty experiments run by the Makata AI Edge Lab against one such
stack---the configuration shipped in \emph{Xylo}, a document-chat app whose 6\,GB tier runs a
single Qwen2.5-1.5B (Q4) model for both summarisation and chat, with EmbeddingGemma-300M for
retrieval, through llama.cpp. Our contribution is not a new algorithm; it is a
\emph{disciplined empirical map} of which levers move answer quality in this regime and which
do not. Three properties make the map trustworthy:

\begin{itemize}
  \item \textbf{Pre-registration.} Each experiment's question, hypotheses and decision rule
  are committed to version control before the first run; deviations are dated, not silent.
  \item \textbf{Published negatives.} Rejected and inconclusive results are written up in the
  same detail as adopted ones---most of the twenty are negative or replication results, which
  is the point.
  \item \textbf{Reproducibility.} Raw per-item outputs, scoring scripts, and the commit hash
  for each registration and write-up are public (Section~\ref{sec:repro}).
\end{itemize}

Our headline findings: a strong embedder dominates retrieval; a cluster of ``hard'' legal
categories is a query-vocabulary problem solvable without any generation; and for clause-style
answers the 1.5B generator---not retrieval---is the limit, so removing it in favour of
extraction helps.

\section{Background and related work}
We build on and test claims from: BM25 and hybrid lexical/dense retrieval
\citep{robertson2009bm25}; contextual retrieval, which prepends model-generated per-chunk
context before indexing \citep{anthropic2024contextual}; structure/clause-aware legal RAG
\citep{strucchunk}; question decomposition \citep{resin}; the sensitivity of small LLMs to
distracting retrieved context \citep{cuconasu2024noise,longctx2025rag}; position effects in
long contexts \citep{liu2023lost}; and reported over-refusal of small on-premises models under
legal authority framing \citep{overrefusal2026}. The benchmark is CUAD \citep{hendrycks2021cuad};
the models are Qwen2.5 \citep{qwen2025} and the Gemma family of open models and embedders
\citep{gemma,embeddinggemma}; inference is via llama.cpp \citep{llamacpp}. Answer quality on
clause-style questions is measured with SQuAD-style token F1 \citep{rajpurkar2016squad}.
Where a technique's source did not evaluate CUAD, EmbeddingGemma, recall@3, or a 1.5B generator,
we treat transfer as unverified and test it.

\section{Setup and method}\label{sec:method}
\paragraph{Stack.} Qwen2.5-1.5B-Instruct Q4\_K\_M (the shipped file), EmbeddingGemma-300M for
dense retrieval, hybrid scoring, top-3 chunks, a 2{,}048-token context. All runs are on a Mac
host (Apple M4 Pro, 24\,GB) through llama.cpp with Metal---the same build the app ships; device
(iPhone) figures are explicitly marked as follow-ups.

\paragraph{Benchmark.} CUAD v1 \citep{hendrycks2021cuad}: 397 commercial contracts, 1{,}325
answerable question--span items. A 464-item subset carries a machine-parsable key fact (a date,
duration, or jurisdiction); the remaining ``hard'' categories (e.g.\ Exclusivity, Non-Compete,
Change Of Control) have clause-text answers with no extractable fact.

\paragraph{Metrics.} \emph{recall@3}: the top-3 chunks cover $\geq$50\% of the gold span.
\emph{key-fact accuracy}: the answer contains the gold fact (the 464 scorable items).
\emph{span-F1}: SQuAD-style token overlap against the gold span, used for the non-fact
categories (absolute values are low for long clauses, so we read the delta, not the number).
Confidence intervals are contract-level bootstraps (10{,}000 resamples, seed 0).

\paragraph{Discipline.} Every experiment is pre-registered; every reported number is tagged
MEASURED (recorded in a run here), OBSERVED (seen, not systematically measured), or INFERRED
(follows from measurements). Verdicts are \emph{adopt}, \emph{reject}, or \emph{inconclusive}
against the registered decision rule.

\section{Retrieval: what helps, what does not}\label{sec:retrieval}
The largest single lever is embedder quality. EmbeddingGemma raises recall@3 from 35\% (the
NLEmbedding fallback) to 61\% (\poc{007}, $+$26 pts, 95\% CI $+$24 to $+$29), at 8.9\,ms/chunk.
Against that, a series of plausible additions do \emph{not} help on this benchmark
(Table~\ref{tab:retrieval}): BM25-only retrieval drops recall 15 points (\poc{010}); having
Qwen2.5-1.5B write a per-chunk context before embedding---Anthropic-style contextual retrieval
with an on-device generator---is a net regression ($-$10.5 pts, \poc{011}); five retrieved
chunks instead of three do not raise answer accuracy (\poc{009}); and clause-aware chunking
gives only $+$2.5 pts (CIs overlap) while \emph{reducing} span integrity, so even that small
gain is not from the registered mechanism (\poc{014}). A recurring cause unites the negatives:
on a small model, extra or mismatched text \emph{dilutes} rather than helps.

\begin{table}[t]
\centering
\small
\begin{tabular}{l l l}
\toprule
PoC & Mechanism tested & Verdict (headline) \\
\midrule
007 & EmbeddingGemma vs NLEmbedding & \textbf{Adopt}: recall 35\%$\rightarrow$61\% ($+$26) \\
010 & BM25-only (drop the embedder)  & Reject: $-$15 pts recall \\
011 & 1.5B-generated chunk context   & Reject (on-device): $-$10.5 pts \\
009 & 5 retrieved chunks vs 3        & Reject: $+$0.4 key-fact (CI $-$3.7,$+$4.4) \\
014 & Clause-aware chunking          & Inconclusive: $+$2.5 (CIs overlap), mechanism refuted \\
013 & Semantic/BM25 weight sweep     & Weight 0.9/0.1 adopted; hard categories invariant \\
\bottomrule
\end{tabular}
\caption{Retrieval experiments. A strong embedder dominates; the other levers do not move
recall on CUAD.}
\label{tab:retrieval}
\end{table}

\section{The hard categories are a query problem}\label{sec:query}
A cluster of categories (Exclusivity around 7--11\%, Non-Compete, Change Of Control) resisted
every retrieval lever above: weight tuning left Exclusivity near 11\% at every setting
(\poc{013}), depth did not help (\poc{009}), and chunk boundaries did not either---page
chunks already contained the gold span whole 95.8\% of the time (\poc{014}). The common cause
is on the \emph{query} side: the CUAD category questions are long shared boilerplate, so the
embedded query vector is dominated by words every category has in common.

Two deterministic, generation-free query transforms move them (Table~\ref{tab:query}).
Template question decomposition---issuing 2--3 fixed sub-queries and unioning their
results---lifts Exclusivity recall 7\%$\rightarrow$24\% (\poc{016}), the first mechanism in the
series to move a hard category at all. Keyword expansion---adding each category's
discriminative vocabulary to the query---does better and cheaper: $+$4.8 pts overall,
$+$8.8 on the weak-category mean, Exclusivity 7\%$\rightarrow$32\% with a single query
(\poc{015}). A fluent natural-language gloss, by contrast, \emph{hurt} ($-$2.0), isolating the
lever as discriminative \emph{terms}, not phrasing. Crucially, the recall gain converts
end-to-end: key-fact accuracy rises $+$3.2 overall and $+$10.7 on Agreement Date (\poc{015},
Arm~B).

\begin{table}[t]
\centering
\small
\begin{tabular}{l l l}
\toprule
PoC & Query transform & Exclusivity recall@3 \\
\midrule
--  & raw category question (baseline) & 7\% \\
016 & template decomposition          & 24\% ($+$17) \\
015 & keyword expansion               & \textbf{32\%} ($+$25) \\
015 & fluent gloss                    & 14\% (overall $-$2.0; rejected) \\
\bottomrule
\end{tabular}
\caption{Query formulation on the hard categories. Injecting discriminative terms---not
rephrasing, generation, chunking, or scoring---is the lever.}
\label{tab:query}
\end{table}

\section{The answer step: extraction beats generation}\label{sec:answer}
Better retrieval exposed a harder limit. On the non-fact hard categories, keyword expansion's
large recall win converts to only a small answer-quality gain (Exclusivity span-F1
2.8$\rightarrow$5.9) and the hard-category mean is flat (\poc{017})---retrieving the right
clause is necessary but not sufficient. The 1.5B generator is now the bottleneck.

Two answer-step changes address it (Table~\ref{tab:answer}). Removing the generator
entirely---selecting the best $\leq$3-sentence passage from the retrieved chunks by
embedding similarity---raises hard-category span-F1 to 20.7, $+$5.0 over the same-retrieval
Qwen arm, while freeing roughly 1\,GB of RAM (\poc{018}). Keeping the generator but forcing it
to emit a verbatim substring, validated by string match with a first-sentence fallback,
reaches 22.1 (\poc{019})---but the validator rejected 61.7\% of generations, so the extractive
fallback, not constrained generation, carries most of the gain. Both point the same way: for
clause-style answers at this scale, extraction is the robust choice, and the generator-free
variant is the cheaper one to ship.

\begin{table}[t]
\centering
\small
\begin{tabular}{l l l}
\toprule
PoC & Answer step (hard categories) & span-F1 \\
\midrule
017 & Qwen2.5-1.5B free generation (baseline)     & 15.4--16.1 \\
018 & generator-free extraction ($\leq$3 sentences) & \textbf{20.7} ($+$5.0), $\sim$1\,GB freed \\
019 & constrained verbatim + validator            & 22.1 ($+$6.6); 61.7\% fell back to extraction \\
\bottomrule
\end{tabular}
\caption{The answer step. Once retrieval is fixed, a generator-free extractive step beats the
1.5B generator on clause-style answers.}
\label{tab:answer}
\end{table}

\section{Prompting: absence and over-refusal}\label{sec:prompting}
Two prompt-level interventions were tested and largely rejected. A firm ``reply exactly \dots\
if the sections do not contain the answer'' instruction did not move the default on-device
model's absence rate without costing answer accuracy (\poc{008}); it helped only the model that
already abstained most. And a reported over-refusal effect---legal ``authority'' framing making
small models refuse more---did \emph{not} transfer: on Qwen2.5-1.5B the authority prompt refused
\emph{least} (26.3\% vs 29.9\% plain vs 35.2\% task-only) and key-fact accuracy was flat
(\poc{020}). A separate comparison of answer backends found the on-device Gemma variant more
willing to say ``absent'' than the platform foundation model (\poc{006}). These are useful
negatives: they rule out cheap prompt-only fixes.

\section{Summary generation}\label{sec:summary}
A sub-thread (\poc{003}--\poc{005}) concerns the app's four-section summary. The quantized
model's tendency to stop after two sections is a \emph{sampler} effect, not a token-budget or
prompt-example effect: the repeat/presence penalties that stop chat rambling also suppress
later section headings, and a summary-specific sampler plus heading-insertion raises complete
four-section output from 39\% to 93\% (\poc{005}). KV-cache quantization, by contrast, was not
the memory constraint it was assumed to be (\poc{001}).

\section{Discussion: a recurring principle}
Across retrieval (\poc{011}, \poc{016}'s cross-reference arm), grounding, and answering, the
same principle recurs: \emph{for a small model, wrong or extra context is a net negative.}
Adding generated per-chunk context, appending resolved cross-references, or injecting
mislabelled entity hints all dilute rather than help; the wins come from \emph{more precise}
inputs (discriminative query terms) and from \emph{removing} an unreliable step (the generator)
rather than adding capability. This is the opposite of the ``add more context'' intuition that
serves large-model RAG.

\section{Threats to validity}
Results are host-stage (Mac, llama.cpp/Metal), not on-device; RAM and latency claims marked
INFERRED need an iPhone run. The benchmark is English US commercial contracts (CUAD); other
domains and languages may behave differently. span-F1 is a proxy for clause-answer quality, read
as a delta only. Several experiments use a single sampler seed in the verdict; where a second
seed was run it is reported. Some transforms use hand-authored per-category templates that may
overfit CUAD phrasing.

\section{Reproducibility and data availability}\label{sec:repro}
All code (Apache-2.0), write-ups and raw results (CC~BY~4.0) are in the \repo{} repository.
Each experiment folder contains its pre-registration, run scripts, raw per-item outputs, and a
write-up citing the commit hash of its registration and results. A versioned snapshot of the
repository is archived on Zenodo (DOI in \texttt{CITATION.cff}). No model weights or datasets
are redistributed; each experiment names its models and the CUAD release it uses.

\section{Conclusion}
On a 6\,GB phone, the levers that help document RAG are a strong embedder, precise query terms
for hard categories, and---counter-intuitively---\emph{removing} the small generator in favour
of extraction for clause-style answers. Many cloud-RAG techniques did not transfer; we report
those negatives in full. We hope the pre-registered, reproducible format is as useful as the
findings.

\section*{Acknowledgements}
Experiments and write-ups were assisted by Claude (Anthropic), disclosed per the lab's
contributing norm; all results are reproducible from the public repository.

\bibliographystyle{plainnat}
\bibliography{references}

\end{document}
